\section{Тестирование и бенчмарк}

Корректность программы проверялась на примере из условия и на сгенерированных наборах данных.
Измеряемая программа сравнивает количество найденных вхождений Z-поиска с наивным перебором на тех
тестах, где наивный алгоритм включён. Во всех таких случаях количество ответов совпало.

Измерения выполнялись по 5 запусков на каждый набор данных. В таблицах приведены медианные значения
в микросекундах.

\subsection{Сравнение с наивным поиском}

На случайных данных наивный алгоритм часто быстро останавливает сравнение уже на первом несовпавшем
слове, поэтому на малых случайных тестах он может быть быстрее по константам. Такой результат не
противоречит асимптотике: худший случай для наивного поиска возникает, когда длинный префикс образца
часто совпадает с текстом.

\begin{center}
\begin{tabular}{|r|r|r|r|r|}
\hline
\rowcolor{lightgray}
Образец & Текст & Z, мкс & Наивн., мкс & Н/Z \\
\hline
1 & 20\,000 & 73 & 25 & \(0.34\times\) \\
\hline
3 & 20\,000 & 40 & 16 & \(0.40\times\) \\
\hline
5 & 20\,000 & 46 & 18 & \(0.39\times\) \\
\hline
10 & 20\,000 & 51 & 24 & \(0.47\times\) \\
\hline
\end{tabular}
\end{center}

Для демонстрации худшего случая был использован текст из одинаковых слов и образец из таких же слов.
В этом случае наивный алгоритм почти в каждой позиции сравнивает весь образец, а Z-функция по-прежнему
строится линейно.

\begin{center}
\begin{tabular}{|r|r|r|r|r|r|}
\hline
\rowcolor{lightgray}
Образец & Текст & Ответов & Z, мкс & Наивн., мкс & Н/Z \\
\hline
10 & 50\,000 & 49\,991 & 199 & 1276 & \(6.41\times\) \\
\hline
25 & 50\,000 & 49\,976 & 206 & 3350 & \(16.26\times\) \\
\hline
50 & 50\,000 & 49\,951 & 228 & 6288 & \(27.58\times\) \\
\hline
\end{tabular}
\end{center}

\subsection{Масштабирование по длине текста}

В следующей серии длина образца фиксирована и равна 5 словам. Менялась только длина текста. Количество
вставленных вхождений составляло примерно одно вхождение на тысячу слов текста.

\begin{center}
\begin{tabular}{|r|r|r|r|}
\hline
\rowcolor{lightgray}
Слов в тексте & Вхождений & Z-функция, мкс \\
\hline
50\,000 & 50 & 59 \\
\hline
100\,000 & 100 & 117 \\
\hline
500\,000 & 500 & 602 \\
\hline
1\,000\,000 & 1000 & 1252 \\
\hline
\end{tabular}
\end{center}

Рост времени близок к линейному: при увеличении текста с \(100000\) до \(1000000\) слов
время выросло примерно в \(10.5\) раза.

\subsection{Масштабирование по длине образца}

Здесь длина текста фиксирована и равна \(200000\) словам, а количество вставленных вхождений равно
\(500\). Менялась длина образца.

\begin{center}
\begin{tabular}{|r|r|r|r|}
\hline
\rowcolor{lightgray}
Длина образца & Слов в тексте & Z-функция, мкс \\
\hline
1 & 200\,000 & 238 \\
\hline
5 & 200\,000 & 250 \\
\hline
10 & 200\,000 & 261 \\
\hline
25 & 200\,000 & 307 \\
\hline
50 & 200\,000 & 380 \\
\hline
\end{tabular}
\end{center}

Увеличение длины образца добавляет слова в начало объединённой последовательности и влияет на число
сравнений при расширении Z-блоков. При этом основной вклад всё равно задаётся длиной текста, поэтому
время растёт умеренно.

\subsection{Перекрывающиеся вхождения}

Отдельно проверялся случай с большим числом перекрывающихся ответов: образец \(x^m\), текст \(x^n\).
В таком тесте вхождение начинается почти в каждой позиции текста.

\begin{center}
\begin{tabular}{|r|r|r|r|}
\hline
\rowcolor{lightgray}
Длина образца & Вхождений & Z-функция, мкс\\
\hline
5 & 99\,996 & 315 \\
\hline
10 & 99\,991 & 335 \\
\hline
20 & 99\,981 & 348 \\
\hline
\end{tabular}
\end{center}

Время остаётся близким для разных длин образца, поскольку текст имеет одинаковую длину, а количество
ответов отличается незначительно. Это соответствует линейной оценке работы Z-функции.
\pagebreak
